% Generated draft: requires mathematical and source review.
\documentclass[11pt]{article}
\usepackage[margin=25mm]{geometry}
\usepackage{fontspec}
\setmainfont{Latin Modern Roman}
\usepackage{amsmath,amssymb,mathtools,unicode-math}
\setmathfont{Latin Modern Math}
\usepackage{enumitem,needspace,xurl,hyperref}
\setlength{\parindent}{0pt}
\setlength{\parskip}{5pt}
\setlength{\emergencystretch}{3em}
\newcommand{\N}{\mathbb{N}}
\newcommand{\Z}{\mathbb{Z}}
\newcommand{\Q}{\mathbb{Q}}
\newcommand{\R}{\mathbb{R}}
\newcommand{\C}{\mathbb{C}}
\newcommand{\F}{\mathbb{F}}
\newcommand{\A}{\mathbb{A}}
\newcommand{\PP}{\mathbb{P}}
\newcommand{\E}{\mathbb{E}}
\newcommand{\eps}{\varepsilon}
\newcommand{\dd}{\,\mathrm{d}}
\begin{document}
\section*{UnsolvedMath ID: 30000963}
\textbf{Source problem code:} catalogue entry
\subsection{Definitions and mathematical statement}
A planar point process may be represented by a random locally finite counting measure $\Xi$ on $\C\cong\R^2$. It is stationary if its law is invariant under every translation. For a nonzero Gaussian random entire function $f$, its zero process is
\[
\Xi_f=\sum_{z:f(z)=0}\operatorname{ord}_z(f)\,\delta_z.
\]
The covariance of $f$ must be specified to identify its law; stationarity of zeros is not automatic for an arbitrary Gaussian entire function. A homogeneous Poisson process of intensity $\rho$ has independent counts on disjoint bounded sets, with each count Poisson of mean $\rho$ times area. A Coulomb-gas comparison additionally requires its interaction, density, inverse temperature and finite-volume or limiting convention.

The retained source asks which stationary planar process is ``most remarkable mathematically'' and how the chaotic analytic zero process should be characterized relative to Poisson and Coulomb-gas models. No ranking, objective criterion, fixed Gaussian covariance or precise characterization property is supplied.

Crucially, the verification field says this sentence is a rhetorical opening to the report's exposition, not a discrete research problem, and promotes no clean mathematical statement. It is therefore recorded as an unrecoverable theorem formulation. The definitions above identify its subject matter but do not manufacture an extremal problem, characterization theorem or uniqueness conjecture from the rhetorical comparison.
\subsection{Short English statement}
This is a rhetorical comparison of stationary planar zero processes with Poisson and Coulomb-gas models. The source explicitly says it is not a discrete research problem, and no objective meaning of ``most remarkable'' or exact characterization is supplied.
\subsection{Sources}
\begin{itemize}
\item Supplied report: \emph{Stochastic Analysis} (2009), Oberwolfach Reports, DOI 10.4171/owr/2008/25, pp. 1361--1424; proposed year 2008. \url{https://doi.org/10.4171/owr/2008/25}.
\item Supplied Sodin--Tsirelson 2002 reference: \url{https://arxiv.org/abs/math/0210090}.
\item Supplied Sodin--Tsirelson 2006 reference: \url{https://doi.org/10.1007/BF02771978}.
\item Local UnsolvedMath registry, permanent ID 30000963, catalogue entry; created 25 August 2026, updated 16 September 2026. The statement status is unrecoverable because the official report uses the sentence rhetorically; the original remains for provenance and no clean statement is promoted. The supplied 21 August 2026 triage likewise says no uniquely defined theorem or objective ranking is specified, despite an open registry label. These qualifications are retained, not converted to an open mathematical conjecture.
\end{itemize}
\end{document}
